\section{Summary and Looking Ahead}
\label{sec:differential-calculus-summary}

Functions describe dependence between quantities, while limits make local behavior mathematically precise. The derivative measures an instantaneous rate of change, and repeated differentiation produces acceleration and higher-order rates. Partial derivatives extend this idea to functions of several variables, whereas the total derivative follows all indirect dependencies along a path.

The operator \(\nabla\) organizes spatial partial derivatives. Acting on a scalar field, it produces the gradient, which points toward the greatest local increase and is perpendicular to level surfaces. The directional derivative then extracts the rate of change along any chosen unit direction.

\begin{booksummary}[Chapter 2 in One Line]
Differential calculus turns continuous change into a local mathematical object that can be interpreted geometrically and physically.
\end{booksummary}

The next chapter applies differential operators to vector fields and connects local derivatives with global integral relations.
